\documentclass[../DoAn.tex]{subfiles}
\begin{document}

Phụ lục này đặc tả chi tiết hai use case bổ sung không trình bày trong Chương~2:
chơi mạng LAN nhiều người (UC2) và xem kết quả backtest (UC5).

\section{Đặc tả UC2 -- Chơi mạng LAN nhiều người}

\begin{table}[htbp]
\centering
\caption{Đặc tả UC2 -- Chơi mạng LAN}
\label{tab:uc2}
\small
\begin{tabular}{lp{10cm}}
\hline
\textbf{Tên use case}   & Chơi mạng LAN nhiều người \\
\textbf{Tác nhân}       & Người chơi chủ phòng (host), người chơi tham gia (client) \\
\textbf{Tiền điều kiện} & Các máy cùng mạng LAN; một máy tạo phòng, các máy khác tham gia \\
\textbf{Luồng chính}    &
1. Host tạo phòng, hệ thống khởi tạo máy chủ Fish-Net và hiển thị mã/địa chỉ phòng. \newline
2. Các client nhập địa chỉ và tham gia phòng. \newline
3. Mỗi người chơi chọn xe tăng và bấm sẵn sàng (ready). \newline
4. Khi tất cả sẵn sàng, host bấm bắt đầu; hệ thống đồng bộ bản đồ và spawn xe tăng cho mọi người chơi. \newline
5. Số agent AI được tính theo công thức 6 agent mỗi người chơi; các agent tấn công Eagle Base chung. \newline
6. Trạng thái xe tăng, đạn và Eagle được đồng bộ giữa các máy theo thời gian thực. \\
\textbf{Luồng phát sinh} &
2a. Client nhập sai địa chỉ → báo lỗi kết nối, quay lại màn hình nhập. \newline
4a. Một client mất kết nối → host nhận thông báo và tiếp tục ván chơi với số người còn lại. \\
\textbf{Hậu điều kiện}  & Ván chơi kết thúc đồng bộ trên mọi máy; hiển thị kết quả chung \\
\hline
\end{tabular}
\end{table}

\section{Đặc tả UC5 -- Xem kết quả backtest}

\begin{table}[htbp]
\centering
\caption{Đặc tả UC5 -- Xem kết quả backtest}
\label{tab:uc5}
\small
\begin{tabular}{lp{10cm}}
\hline
\textbf{Tên use case}   & Xem kết quả backtest \\
\textbf{Tác nhân}       & Người dùng (nghiên cứu viên) \\
\textbf{Tiền điều kiện} & Đã chạy ít nhất một đợt backtest và sinh tệp kết quả \\
\textbf{Luồng chính}    &
1. Người dùng mở thư mục \texttt{BacktestResults/}. \newline
2. Mở tệp \texttt{backtest\_summary\_*.csv} để xem kết quả tổng hợp theo từng run. \newline
3. Mở tệp \texttt{backtest\_chart\_*.html} để xem biểu đồ trực quan trên trình duyệt. \newline
4. Mở tệp \texttt{backtest\_agents\_*.csv} để xem chi tiết chỉ số từng agent. \\
\textbf{Luồng phát sinh} &
3a. Nếu môi trường thiếu Python/matplotlib, tệp PNG không sinh được nhưng HTML vẫn hiển thị bảng số liệu. \\
\textbf{Hậu điều kiện}  & Người dùng nắm được kết quả định lượng của đợt thực nghiệm \\
\hline
\end{tabular}
\end{table}

\end{document}
